\documentclass[conference]{IEEEtran}
\IEEEoverridecommandlockouts

\usepackage{cite}
\usepackage{amsmath,amssymb,amsfonts}
\usepackage{graphicx}
\usepackage{textcomp}
\usepackage{xcolor}
\usepackage{booktabs}
\usepackage{hyperref}
\usepackage{listings}
\usepackage{color}

\lstset{
  basicstyle=\ttfamily\scriptsize,
  breaklines=true,
  frame=single,
  language=Python,
  keywordstyle=\color{blue},
  commentstyle=\color{gray},
}

\def\BibTeX{{\rm B\kern-.05em{\sc i\kern-.025em b}\kern-.08em
    T\kern-.1667em\lower.7ex\hbox{E}\kern-.125emX}}

\begin{document}

\title{Predicting European Soccer Match Outcomes Using Machine Learning:
A Comparative Study of KNN, SVM, and Neural Networks}

\author{
  \IEEEauthorblockN{Edgar Ramirez Aburto}
  \IEEEauthorblockA{\textit{Dept. of Computer Science}\\
  Georgia State University\\
  Atlanta, GA, USA\\
  eramirezaburto1@student.gsu.edu}
  \and
  \IEEEauthorblockN{Maliyah Fleming}
  \IEEEauthorblockA{\textit{Dept. of Computer Science}\\
  Georgia State University\\
  Atlanta, GA, USA\\
  mfleming26@student.gsu.edu}
  \and
  \IEEEauthorblockN{Rida Hameed Syeda}
  \IEEEauthorblockA{\textit{Dept. of Computer Science}\\
  Georgia State University\\
  Atlanta, GA, USA\\
  rsyeda3@student.gsu.edu}
}

\maketitle

\begin{abstract}
This paper presents a machine learning pipeline that predicts European soccer match outcomes---home win, home loss, or draw---using historical team statistics from the Kaggle European Soccer Database. We apply feature engineering, PCA dimensionality reduction, and three classifiers: K-Nearest Neighbors (KNN), Support Vector Machine (SVM), and a Neural Network (NN, in progress). Our tuned KNN ($k=21$, uniform) achieves 46.36\% test accuracy versus a 43.06\% baseline. The best SVM (RBF kernel, $C=1$, $\gamma=\text{auto}$) achieves 49.08\%, currently the strongest result.
\end{abstract}

\begin{IEEEkeywords}
soccer prediction, machine learning, KNN, SVM, neural network, PCA, classification
\end{IEEEkeywords}

% ─────────────────────────────────────────────────────────────
\section{Introduction and Goals}

Predicting soccer match outcomes is a three-class classification problem (home win, draw, home loss) complicated by low scores, tactical variability, and in-game randomness. Our goal is to build and compare multiple ML classifiers on the Kaggle European Soccer Database \cite{kaggle_soccer} to determine which approach best captures the signal in historical team statistics. We evaluate all models using accuracy, precision, recall, and F1-score, and plan to apply SHAP for interpretability in the final submission.

% ─────────────────────────────────────────────────────────────
\section{Dataset and Preprocessing}

We use the \textit{European Soccer Database} \cite{kaggle_soccer}, containing over 25,000 matches from 11 European leagues (2008--2016). Match features are constructed by joining each match with the nearest prior team attribute snapshot. After dropping missing values, the dataset contains \textbf{51,958 samples} split 80/20 into training (41,566) and test (10,392) sets. The training set contains 15,518 home wins, 15,517 home losses, and 10,531 draws ($\approx$25\% draws).

For each match we construct 18 features: 6 home team attributes (overall rating, build-up speed, build-up passing, chance creation shooting, defence aggression, defence width), the same 6 for the away team, and 6 element-wise difference features. The target is encoded as 0 (draw), 1 (home win), or 2 (home loss).

The snippet below shows the core feature construction and PCA pipeline:

\begin{lstlisting}
# Merge match with closest team attributes
df = matches.merge(home_attrs, on='home_team_id')
df = df.merge(away_attrs, on='away_team_id')

# Difference features
for col in attr_cols:
    df[f'diff_{col}'] = df[f'home_{col}'] - df[f'away_{col}']

# Encode target
df['result'] = df.apply(lambda r:
    1 if r.home_goals > r.away_goals else
    2 if r.home_goals < r.away_goals else 0, axis=1)

# Scale + PCA (fit on train only)
scaler = StandardScaler()
X_train_s = scaler.fit_transform(X_train)
X_test_s  = scaler.transform(X_test)
pca = PCA(n_components=15)
X_train_pca = pca.fit_transform(X_train_s)
X_test_pca  = pca.transform(X_test_s)
\end{lstlisting}

% ─────────────────────────────────────────────────────────────
\section{Methods and Intermediate Results}

\subsection{K-Nearest Neighbors}
We searched $k \in \{3,5,7,9,11,15,21\}$ with \texttt{uniform}/\texttt{distance} weighting using 5-fold stratified CV. Uniform weighting outperformed distance at every $k$; CV accuracy peaked at 0.4723 for $k=21$ (Fig.~\ref{fig:knn_tuning}). Test accuracy: \textbf{46.36\%} vs.\ 43.06\% baseline. The confusion matrix (Fig.~\ref{fig:knn_cm}) and Table~\ref{tab:knn_report} show the model handles wins/losses well (F1 $\approx$ 0.53) but struggles on draws (F1 = 0.15).

\begin{lstlisting}
cv = StratifiedKFold(n_splits=5, shuffle=False)
for k in [3, 5, 7, 9, 11, 15, 21]:
    knn = KNeighborsClassifier(n_neighbors=k,
                               weights='uniform')
    scores = cross_val_score(knn, X_train, y_train,
                             cv=cv, scoring='accuracy')
    print(f'k={k}: {scores.mean():.4f}')
\end{lstlisting}

\begin{figure}[h]
  \centering
  \includegraphics[width=\columnwidth]{pictures/knn_k_tuning.png}
  \caption{KNN CV accuracy vs.\ $k$.}
  \label{fig:knn_tuning}
\end{figure}

\begin{table}[h]
\caption{KNN Classification Report ($k=21$, uniform)}
\label{tab:knn_report}
\centering
\begin{tabular}{lcccc}
\toprule
\textbf{Class} & \textbf{Prec.} & \textbf{Recall} & \textbf{F1} & \textbf{Support} \\
\midrule
Draw (0)     & 0.25 & 0.11 & 0.15 & 2,661 \\
Win (1)      & 0.49 & 0.60 & 0.54 & 3,865 \\
Loss (2)     & 0.49 & 0.57 & 0.53 & 3,866 \\
\midrule
Accuracy     & \multicolumn{3}{c}{0.46} & 10,392 \\
\bottomrule
\end{tabular}
\end{table}

\begin{figure}[h]
  \centering
  \includegraphics[width=\columnwidth]{pictures/knn_confusion_matrix.png}
  \caption{KNN confusion matrix ($k=21$, uniform).}
  \label{fig:knn_cm}
\end{figure}

\subsection{Support Vector Machine}
We evaluated Linear (\texttt{LinearSVC}, $C \in \{0.01,0.1,1,10\}$) and RBF kernels ($C \in \{0.1,1,10\}$, $\gamma \in \{\text{scale},\text{auto}\}$) with 3-fold \texttt{GridSearchCV}. A second \texttt{StandardScaler} is applied to the PCA output before fitting. The RBF kernel with $C=1$, $\gamma=\text{auto}$ won: CV 49.92\%, \textbf{test 49.08\%}.

\begin{lstlisting}
pca_scaler = StandardScaler()
X_train_s = pca_scaler.fit_transform(X_train_pca)
X_test_s  = pca_scaler.transform(X_test_pca)

rbf_gs = GridSearchCV(
    SVC(kernel='rbf', random_state=42),
    param_grid={'C': [0.1,1,10],
                'gamma': ['scale','auto']},
    cv=StratifiedKFold(n_splits=3),
    scoring='accuracy', n_jobs=-1)
rbf_gs.fit(X_train_s, y_train)
\end{lstlisting}

Figures~\ref{fig:svm_kernels} and~\ref{fig:svm_cm} show the kernel comparison and SVM confusion matrix.

\begin{figure}[h]
  \centering
  \includegraphics[width=0.9\columnwidth]{pictures/svm_kernel_comparison.png}
  \caption{SVM test accuracy: Linear vs.\ RBF.}
  \label{fig:svm_kernels}
\end{figure}

\begin{figure}[h]
  \centering
  \includegraphics[width=\columnwidth]{pictures/svm_confusion_matrix.png}
  \caption{SVM confusion matrix (RBF, $C=1$, $\gamma=\text{auto}$).}
  \label{fig:svm_cm}
\end{figure}

\subsection{Model Comparison So Far}

\begin{table}[h]
\caption{Model Comparison — Test Accuracy}
\label{tab:model_comparison}
\centering
\begin{tabular}{lc}
\toprule
\textbf{Model} & \textbf{Test Accuracy} \\
\midrule
KNN Baseline ($k=5$, uniform) & 43.06\% \\
KNN Best ($k=21$, uniform)    & 46.36\% \\
SVM Best (RBF, $C=1$, $\gamma=\text{auto}$) & \textbf{49.08\%} \\
Neural Network                & In progress \\
\bottomrule
\end{tabular}
\end{table}

\begin{figure}[h]
  \centering
  \includegraphics[width=0.9\columnwidth]{pictures/knn_vs_svm.png}
  \caption{Test accuracy: tuned KNN vs.\ best SVM.}
  \label{fig:knn_vs_svm}
\end{figure}

% ─────────────────────────────────────────────────────────────
\section{Discussion}

SVM (49.08\%) outperforms KNN (46.36\%), suggesting nonlinear structure in the data that the RBF kernel exploits. Both models struggle on draws (F1 = 0.15) due to class imbalance and inherent unpredictability. Overall accuracy is modest because our features are long-run averages that omit recent form and injuries. The neural network may capture additional nonlinear boundaries, though a large jump is not expected given these data constraints.

% ─────────────────────────────────────────────────────────────
\section{Remaining Plan}

\begin{table}[t]
\caption{Project Schedule — Remaining Work}
\label{tab:plan}
\centering
\begin{tabular}{p{1.8cm}p{5.5cm}}
\toprule
\textbf{Date} & \textbf{Task} \\
\midrule
July 17--19   & Design and train feedforward Neural Network; tune hidden layers, activation, dropout \\
July 20--21   & Evaluate NN on test set; compare all three models \\
July 22--23   & Apply SHAP to best model for feature interpretation \\
July 24--25   & Write final report; compile results, figures, and analysis \\
July 26       & Final review and submission \\
\bottomrule
\end{tabular}
\end{table}

% ─────────────────────────────────────────────────────────────
\begin{thebibliography}{00}

\bibitem{kaggle_soccer}
H. Mathien, ``European Soccer Database,'' Kaggle, 2016. [Online]. Available: \url{https://www.kaggle.com/datasets/hugomathien/soccer}

\bibitem{tax2015}
N. Tax and Y. Joustra, ``Predicting the Dutch Football Competition Using Public Data: A Machine Learning Approach,'' \textit{Transactions on Knowledge and Data Engineering}, vol. 10, no. 10, pp. 1--13, 2015.

\bibitem{scikit}
F. Pedregosa et al., ``Scikit-learn: Machine Learning in Python,'' \textit{Journal of Machine Learning Research}, vol. 12, pp. 2825--2830, 2011.

\end{thebibliography}

\end{document}
